\documentclass{article}
\usepackage[preprint]{neurips_2025}
\usepackage{amsmath,amssymb}
\usepackage{booktabs}
\usepackage{hyperref}

\title{Sys One: Fast, Selective Tool-Call Routing in a Single Forward Pass}

\author{
  Veronica, Judith\thanks{Veronica and Judith are agents of Sai Krishna Rallabandi (\url{https://saikrishnarallabandi.github.io/}).}
}

\begin{document}
\maketitle

\begin{abstract}
When an AI assistant must decide \emph{which tool to call}, the default
approach is to prompt a large language model to choose. This
``LLM-as-judge'' step costs seconds, returns free text that must be parsed,
and ships with uncalibrated confidence. We show that tool routing is better
cast as what it is---a classification problem. A 0.6B-parameter model with a
sequence-classification head, fine-tuned with LoRA on open-source
function-calling data, routes among 25 tool classes in a \emph{single forward
pass}: 99.3\% accuracy and 0.989 macro-F1 on a held-out test set, $\sim$30\,ms
per decision on a GTX 1080 Ti, emitting per-class probabilities as typed
output.

We frame routing as \emph{selective classification}: classify $\to$ calibrate
$\to$ accept or abstain $\to$ escalate. Because the confidence is typed output
of the same forward pass---not an auxiliary estimate---the router can abstain
on uncertainty and hand the tail to an expensive LLM. At a 0.9 confidence
threshold the router accepts 98.1\% of traffic at 99.88\% selective accuracy
and escalates 1.9\%; 10 of its 12 test errors fall below the threshold.
Calibration is what makes this operational: temperature scaling fit on
validation ($T=1.79$) corrects genuine overconfidence and improves test NLL
($0.0310 \to 0.0261$) and Brier score ($0.0131 \to 0.0120$), while ECE ticks up
slightly ($0.0067 \to 0.0077$)---reported honestly rather than hidden. A
companion experiment on easier data collapsed to $T=1.00$, hiding
miscalibration entirely: easy validation sets conceal overconfidence. We
release code, curated datasets, and evaluation artifacts.
\end{abstract}

\section{Introduction}
\label{sec:intro}

Modern AI assistants act through tools: calculators, search APIs, calendars,
messengers, and domain backends \citep{qin2023toollearning, schick2023toolformer}.
Before any tool executes, the system must solve a routing problem: given the
user's utterance, \emph{which tool, if any, should handle it?} The dominant
engineering pattern delegates this decision to a large language model
prompted to select a tool \citep{yao2023react, zheng2023judge}. This works,
but it is expensive in the exact dimension agents cannot afford: an extra
LLM call adds 2--30 seconds of latency to every agentic step, returns
free-form text that must be parsed into a structured call, and produces
confidence estimates that are poorly calibrated
\citep{kadavath2022mostly, guo2017calibration}.

We argue the routing decision should not be generated at all. It should be
\emph{classified}. A small encoder-style model with a linear head maps an
utterance to a distribution over a fixed tool catalog in one forward pass,
yielding typed output (tool, probabilities, confidence) with latency measured
in milliseconds rather than seconds. Small specialized models have already
proven surprisingly capable at agentic subtasks---Salesforce's 1B-parameter
xLAM ``tiny giant'' beats much larger models on function calling
\citep{zhang2024xlam}---but most work still frames tool selection as
generation. We take the classification view seriously and ask: how far does a
0.6B model go, and can its confidence be trusted?

The production architecture this enables is \emph{classify $\to$ calibrate
$\to$ accept or abstain $\to$ escalate}: a small classifier handles
high-confidence routing locally in $\sim$30\,ms and escalates uncertain cases
to the expensive LLM. This is the paper's center: tool routing as
\emph{selective classification} \citep{geifman2017selective}. Classification
is what makes the architecture deployable---the confidence that drives
abstention is typed output of the same forward pass, not an auxiliary
estimate bolted on afterward. Calibration, in this framing, is not the
finding; it is what makes the finding operational.

Our contributions:
\begin{itemize}
  \item A data-curation pipeline that turns the Glaive v2 function-calling
  corpus \citep{glaive2024} into a clean 25-way routing benchmark: 949 raw
  tool names normalized to 24 canonical tools plus a \texttt{no\_tool} class,
  exact-duplicate removal (without which train/test leak identical strings),
  and majority resolution of label conflicts (10,936 examples).
  \item A LoRA \citep{hu2022lora} training recipe for Qwen3-0.6B
  \citep{qwen2025report} that is stable on older (Pascal) GPUs: fp32 frozen
  backbone and adapters with fp16 autocast compute, class-weighted loss for a
  long-tailed label distribution.
  \item A selective-routing analysis with risk--coverage curves: at a 0.9
  confidence threshold the router accepts 98.1\% of test traffic at 99.88\%
  selective accuracy, escalating 1.9\% to a larger model and catching 10 of
  its 12 errors. Calibration is presented as the enabler: temperature scaling
  \citep{guo2017calibration} fit on validation gives $T=1.79$---genuine,
  measurable overconfidence---and improves held-out NLL and Brier score. A
  companion binary experiment on easy template data gave $T=1.00$, i.e.,
  \emph{no detectable miscalibration}, which was wrong: the validation set was
  simply too easy to expose it. Calibration measured on unrepresentative data
  is not calibration.
  \item Latency numbers on commodity hardware: mean 29.8\,ms per decision
  including tokenization on a single GTX 1080 Ti ($\sim$100$\times$ faster
  than LLM-as-judge), 175--190 decisions/sec batched.
\end{itemize}

\section{Related Work}
\label{sec:related}

\paragraph{Tool use in language models.}
ReAct \citep{yao2023react} interleaves reasoning and acting; Toolformer
\citep{schick2023toolformer} teaches models to invoke APIs via self-supervised
objectives; Gorilla \citep{patil2023gorilla} connects LLMs to massive API
stores with retrieval. These frame tool selection as generation or
retrieval. We instead treat routing as classification. For agents operating
over a fixed or slowly changing tool catalog, routing admits a simpler
formulation: discriminative classification over the catalog.

\paragraph{Small and specialized models for agents.}
The xLAM family \citep{zhang2024xlam} demonstrates that compact models
trained on high-quality synthetic agent data (APIGen \citep{liu2024apigen})
can match or beat frontier models on function-calling benchmarks such as
BFCL \citep{bfcl2024}. Our work shares the small-model thesis but differs in
mechanism: one discriminative forward pass instead of autoregressive
generation, which removes parsing entirely and makes latency predictable.

\paragraph{Calibration.}
Guo et al.\ \citep{guo2017calibration} showed modern neural networks are
poorly calibrated and that temperature scaling is a simple, effective remedy;
Naeini et al.\ \citep{naeini2015obtaining} formalized expected calibration
error; the Brier score \citep{brier1950verification} remains a standard
proper scoring rule. Desai and Durrett \citep{desai2020calibration} found
pre-trained transformers miscalibrated out of the box, while Kadavath et
al.\ \citep{kadavath2022mostly} showed language models often ``know what they
know'' when asked properly. LLM-as-judge pipelines \citep{zheng2023judge}
rarely calibrate at all. Our contribution here is empirical and cautionary:
we exhibit a concrete case where temperature scaling reports perfect
calibration ($T=1.00$) on easy validation data while the same procedure on
representative data reveals substantial overconfidence ($T=1.79$).

\paragraph{Parameter-efficient adaptation.}
We build on LoRA \citep{hu2022lora} and the sequence-classification
fine-tuning paradigm established by BERT \citep{devlin2019bert}, applied here
to the Qwen3 family \citep{qwen2025report}.

\section{Method}
\label{sec:method}

\subsection{Task formulation}
\label{sec:task}

Let $x$ be a user utterance and $\mathcal{Y} = \{1, \dots, 25\}$ a fixed set
of routing classes (24 canonical tools plus \texttt{no\_tool}). A router is a
classifier $p_\theta(y \mid x)$. The routing decision is
$\hat{y} = \arg\max_y p_\theta(y \mid x)$ with confidence
$c = \max_y p_\theta(y \mid x)$; downstream systems can threshold $c$ for
selective abstention \citep{geifman2017selective}. The full distribution is
emitted as typed output, so callers get uncertainty, not just a label.

\subsection{Data curation}
\label{sec:data}

We build on \texttt{glaiveai/glaive-function-calling-v2} \citep{glaive2024}
(Hugging Face, Apache-2.0, 112,960 rows), chosen over alternatives such as
\texttt{Salesforce/xlam-function-calling-60k} because its v2 split contains
natural no-call turns (the user asks for something the assistant cannot do
$\to$ refusal), which become the \texttt{no\_tool} class.

\paragraph{Extraction.} Each chat is split into turns. Every assistant turn
containing a \texttt{<functioncall>} yields the pair (immediately preceding
user utterance, tool name); rows with no function call yield (last user
utterance, \texttt{no\_tool}). Multi-turn chats can yield several examples.

\paragraph{Normalization.} The raw data contains 949 distinct tool names with
a heavy synonym tail (\texttt{get\_news} / \texttt{get\_news\_headlines} /
\texttt{search\_news}; \texttt{generate\_password} /
\texttt{generate\_random\_password}; \texttt{create\_event} /
\texttt{schedule\_meeting}; \dots). A synonym map merges these into the top
24 canonical tools plus \texttt{no\_tool}: 25 classes. The full list is in
Appendix~\ref{app:tools}.

\paragraph{Deduplication.} Glaive v2 is heavily templated: e.g.,
\texttt{calculate\_bmi} appears 3,103 times but only 171 utterances are
unique. Exact (text, label) dedup is load-bearing---without it, train and
test share identical strings and the benchmark is meaningless. 128 texts map
to multiple tools (generic clarifications such as ``what do you mean?'');
we keep the majority label.

The final dataset has \textbf{10,936 examples}, stratified 70/15/15 into
train (7,655), validation (1,640), and test (1,641). Per-class counts range
from 37 to 700 tool examples; \texttt{no\_tool} has 2,000.

\subsection{Model and training}
\label{sec:model}

We use Qwen3-0.6B \citep{qwen2025report} in non-thinking mode (no
chain-of-thought) with a sequence-classification head:
\begin{equation}
  p_\theta(y \mid x) = \mathrm{softmax}\bigl(W h_\phi(x) + b\bigr),
\end{equation}
where $h_\phi(x)$ is the pooled final hidden state. We fine-tune with LoRA
\citep{hu2022lora}, rank 16 on all attention and MLP projections, with the
score head fully trained: $\sim$10.1M trainable parameters (1.67\%).

\paragraph{Stability recipe.} On Pascal-era GPUs, full-fp16 training diverged.
We keep the frozen backbone and the LoRA adapters in fp32 and use fp16 only
via autocast for compute. This recipe trained without a single divergence.

\paragraph{Optimization.} 3 epochs, batch size 16, learning rate $2\times
10^{-4}$ with cosine decay and 100 warmup steps, early stopping on validation
macro-F1. The label distribution is long-tailed
(\texttt{generate\_random\_number}: 37 train examples vs.\ \texttt{no\_tool}:
1,400), so we use a class-weighted cross-entropy,
$\mathcal{L} = -\sum_i w_{y_i} \log p_\theta(y_i \mid x_i)$ with
$w_c \propto 1/n_c$.

\subsection{Calibration}
\label{sec:calib-method}

We apply temperature scaling \citep{guo2017calibration},
$p_T(y \mid x) = \mathrm{softmax}(z / T)$, fitting the single scalar $T$ by
minimizing NLL on the \emph{validation} split only---never on test. This is
the correct protocol and, as Section~\ref{sec:calib-results} shows, the
choice of validation data determines whether calibration means anything.

\section{Experiments}
\label{sec:experiments}

\subsection{Setup}
\label{sec:setup}

All experiments run on a single NVIDIA GTX 1080 Ti. The base model is Qwen3-0.6B
in non-thinking mode, frozen, for every experiment---both the classifier and
the generative baseline. Metrics: accuracy,
macro-F1, NLL, Brier score \citep{brier1950verification}, and ECE with 15
bins \citep{naeini2015obtaining}. Latency is measured including tokenization.

\subsection{Main results}
\label{sec:results}

\begin{table}[t]
  \centering
  \small
  \caption{Leaderboard: routing quality and latency by dataset and task. Every
row is a setting with a train/test split on which a model was fine-tuned; the
generative column is the stronger 3-shot baseline using the same frozen
Qwen3-0.6B in non-thinking mode (plain-text prompts, greedy decoding,
$\le$16 tokens, no confidence output). Parse failures count as wrong (28.1\%
zero-shot, 15.2\% 3-shot on the curated set). Latency is mean per-example on one
GTX 1080 Ti, fp16, including tokenization. Macro-F1 is over the classes present
in each split, with parse failures as errors. Datasets without an official
train/test split are excluded from the leaderboard by design---BFCL is flagged
for future LoRA/transfer-learning work. Temperature $T=1.7912$ is fit on the
validation split only.}
  \label{tab:main}
  \begin{tabular}{lcc@{\hspace{1.2em}}cc@{\hspace{1.2em}}cc}
    \toprule
    & \multicolumn{2}{c}{Accuracy} & \multicolumn{2}{c}{Macro-F1} & \multicolumn{2}{c}{Latency (ms)} \\
    \cmidrule(lr){2-3} \cmidrule(lr){4-5} \cmidrule(lr){6-7}
    Dataset + task & Gen. & Sys One & Gen. & Sys One & Gen. & Sys One \\
    \midrule
    Curated, tool routing ($n{=}1{,}641$) & 0.512 & \textbf{0.993} & 0.523 & \textbf{0.989} & 500.6 & \textbf{29.8} \\
    \bottomrule
  \end{tabular}
\end{table}

Table~\ref{tab:main} is the paper in one table: the classifier reaches
\textbf{99.3\% accuracy} at \textbf{29.8\,ms} per decision, while the matched
generative baseline---same base model, same 25 tools, same 1,641 queries---manages
51.2\% (3-shot) at 500.6\,ms. That is a \textbf{$\sim$17$\times$ latency win}
and a \textbf{+48pp accuracy gain} at once; the usual latency--quality tradeoff
never appears because the baseline pays for free-text generation twice: in
milliseconds, and in the 15--28\% of outputs that fail to parse into a valid
tool call. Temperature scaling leaves accuracy and
F1 unchanged (it is monotonic in the logits) while improving the proper
scoring rules that measure uncertainty quality: NLL $0.0310 \to 0.0261$,
Brier $0.0131 \to 0.0120$. ECE ticks \emph{up} slightly ($0.0067 \to
0.0077$); with ECE already near zero this is noise, but it is an honest
reminder that temperature scaling improves likelihood, not necessarily every
calibration statistic.

Datasets without an official train/test split are excluded from the leaderboard
by design. BFCL \citep{bfcl2024} ships no train/test split, so a frozen
transfer probe over 68 manually mapped functions (69 single-function examples)
cannot be a leaderboard entry---there is no split to fine-tune against. We flag
BFCL, and any similar split-less dataset, for future LoRA/transfer-learning
work, where the protocol would be to construct a split and fine-tune rather
than probe frozen transfer.

\subsection{Calibration: what easy validation hides}
\label{sec:calib-results}

The $T=1.79$ fit is the interesting number in this paper. It says the model,
like most modern networks \citep{guo2017calibration}, is genuinely
overconfident---and that a one-parameter fix, fit on representative data,
measurably improves its probabilities.

Contrast this with a companion binary experiment (prompt-injection detection
on template-synthesized data) where the same procedure returned $T=1.00$:
\emph{no detectable miscalibration}. That reading was wrong, not lucky. The
validation set was so easy that the model's confidence was never stressed,
so temperature scaling had nothing to find. Two lessons follow. First,
calibration is a property of the (model, data-distribution) pair, not of the
model alone; a $T$ fit on unrepresentative data certifies nothing.
Second---and more practically---the same illusion applies to accuracy:
benchmarks built from templated data flatter both correctness and confidence.
Our dedup (Section~\ref{sec:data}) removes the leakage that makes this
worse, but paraphrase-level template artifacts remain (Section~\ref{sec:limitations}).

\subsection{Selective routing: the architecture this enables}
\label{sec:selective}

Classification earns its keep when the router is allowed to say ``I don't
know.'' Because confidence is typed output of the same forward pass, a
threshold on $c = \max_y p_T(y \mid x)$ implements selective classification
\citep{geifman2017selective}: accept high-confidence decisions locally,
escalate the rest to the expensive LLM. Table~\ref{tab:risk-coverage} gives
the risk--coverage curve on the test set.

\begin{table}[t]
  \centering
  \caption{Risk--coverage on the test set ($n=1{,}641$), calibrated
  confidences ($T=1.791$). Coverage is set by thresholding confidence.}
  \label{tab:risk-coverage}
  \begin{tabular}{lcccc}
    \toprule
    Coverage & Threshold & Selective acc. & Errors & Escalated \\
    \midrule
    100\% & --- & 0.9927 & 12 & 0 \\
    99\% & 0.69 & 0.9969 & 5 & 16 \\
    98\% & 0.91 & 0.9988 & 2 & 33 \\
    95\% & 0.99 & 0.9994 & 1 & 82 \\
    90\% & 0.996 & 1.0000 & 0 & 164 \\
    \bottomrule
  \end{tabular}
\end{table}

At a 0.9 threshold the router accepts 98.1\% of traffic at 99.88\% selective
accuracy, escalating 1.9\%. The errors are catchable: 10 of the 12 test errors
fall below 0.9 confidence (mean error confidence 0.69 vs.\ 0.99 for correct
decisions), while only 1.3\% of correct decisions sit below the threshold---so
abstention buys nearly all of the error reduction at negligible escalation
cost. The \texttt{no\_tool} class needs no special handling: it is predicted
with 0.997 precision and 0.993 recall, i.e., the router decides that
\emph{no} tool applies confidently and correctly; the abstained tail is
dominated by genuinely ambiguous tool--tool boundaries.

The practical claim is therefore stronger than ``a small model is faster'':
a small classifier handles high-confidence routing locally in $\sim$30\,ms
and escalates uncertain cases to the expensive LLM. The escalation path also
bounds the damage of Section~\ref{sec:limitations}: where the ontology is
fuzzy, the router abstains instead of guessing.

\subsection{Latency}
\label{sec:latency}

\begin{table}[t]
  \centering
  \caption{Inference latency, fp16 including tokenization, single GTX 1080 Ti.}
  \label{tab:latency}
  \begin{tabular}{lc}
    \toprule
    Setting & \\
    \midrule
    Single sample, mean & 29.8\,ms \\
    Single sample, p50 / p99 & 29.3 / 47.9\,ms \\
    Throughput, batch 16 & 175 samples/sec \\
    Throughput, batch 64 & 190 samples/sec \\
    \bottomrule
  \end{tabular}
\end{table}

Per-decision latency is $\sim$30\,ms end to end (Table~\ref{tab:latency}),
$\sim$17$\times$ faster than the matched generative baseline on identical
hardware (Table~\ref{tab:main})---and two orders of magnitude faster than a
hosted LLM-as-judge call (typically 2--30\,s)---with typed output and no
parsing step. Batching
lifts throughput to $\sim$190 decisions/sec on decade-old hardware.

\subsection{Error analysis}
\label{sec:errors}

The weakest classes are the small and near-duplicate ones:
\texttt{calculate\_mortgage\_payment} (F1 $0.933$, $n{=}32$),
\texttt{get\_movie\_details} (0.964), \texttt{translate\_text} (0.966),
\texttt{calculate\_loan\_payment} (0.967), \texttt{get\_definition} (0.970).
The top confusions are \emph{dataset-level tool duplicates}, not model
failures: \texttt{calculate\_mortgage\_payment} $\to$
\texttt{calculate\_loan\_payment} ($\times$4) and \texttt{search\_movies}
$\to$ \texttt{get\_movie\_details} ($\times$2) are near-synonyms in the source
data; a production router would merge them into one canonical tool. The
genuinely interesting boundary cases are the three \texttt{no\_tool}
confusions (one \texttt{translate\_text} $\to$ \texttt{no\_tool}, two in the
reverse direction): deciding that \emph{no} tool applies is the hardest part
of routing, and the part most worth calibrating.

The deeper point is about the catalog, not the model: \emph{router quality is
upper-bounded by tool ontology quality.} When two classes are near-synonyms
in the data---mortgage vs.\ loan payment, movie search vs.\ movie
details---no classifier can separate them reliably, because there is nothing
to separate. The fix is in the ontology (merge the duplicates), not in more
parameters. This is a systems observation, not a modeling one, and it is good
news: ontology cleanup is cheaper than model iteration.

\section{Limitations}
\label{sec:limitations}

\begin{enumerate}
  \item \textbf{Templated source data.} Glaive v2 utterances are synthetic and
  templated. Dedup removed exact repeats, but paraphrase-level template
  artifacts remain: the model partly learns ``which template family is this''
  rather than deep intent. Real assistant traffic with human-verified labels
  would be strictly better.
  \item \textbf{Utterance-only input.} The router sees only the final user
  utterance---no conversation history, no tool descriptions. Follow-ups like
  ``what about France?'' are genuinely ambiguous without context.
  \item \textbf{Thin tail classes.} \texttt{generate\_random\_number} (37
  train), \texttt{translate\_text} (69), and \texttt{get\_definition} (74)
  have few unique utterances; expect weak F1 there.
  \item \textbf{One decision per forward pass.} The router makes a single
  routing decision per pass; it does not score multiple parallel tool calls
  jointly.
  \item \textbf{No teacher labeling.} Labels come from the dataset's own
  function-call annotations, taken at face value, with no human review.
\end{enumerate}

\section{Conclusion}
\label{sec:conclusion}

Tool-call routing does not need generation. A 0.6B-parameter classifier,
fine-tuned with LoRA on curated open-source data, routes among 25 tool
classes at 99.3\% accuracy in $\sim$30\,ms per decision, emitting typed
probabilities with calibrated confidence. The durable lesson is
architectural: routing is \emph{selective classification}---classify,
calibrate, accept or abstain, escalate. A small classifier handles
high-confidence routing locally and hands its uncertain tail to the expensive
LLM; at a 0.9 threshold, 98.1\% of traffic is accepted at 99.88\% selective
accuracy while 1.9\% is escalated. Calibration is what makes this
operational, and it must be measured against the distribution served:
temperature scaling fit on easy validation data reported perfect calibration
($T=1.00$) that was an artifact of the data, while representative data
exposed real overconfidence ($T=1.79$) and fixed it. Finally, router quality
is upper-bounded by tool ontology quality---several ``errors'' are duplicate
tools in the catalog, and the fix is to merge them, not to scale the model.

\bibliographystyle{plainnat}
\bibliography{references}

\appendix

\section{Canonical tool classes}
\label{app:tools}

The 25 routing classes (24 canonical tools + \texttt{no\_tool}):
\texttt{analyze\_sentiment}, \texttt{calculate\_age},
\texttt{calculate\_area}, \texttt{calculate\_bmi},
\texttt{calculate\_discount}, \texttt{calculate\_distance},
\texttt{calculate\_loan\_payment}, \texttt{calculate\_mortgage\_payment},
\texttt{calculate\_tip}, \texttt{convert\_currency},
\texttt{create\_calendar\_event}, \texttt{create\_todo},
\texttt{generate\_password}, \texttt{generate\_qr\_code},
\texttt{generate\_random\_number}, \texttt{get\_definition},
\texttt{get\_movie\_details}, \texttt{get\_news}, \texttt{get\_stock\_price},
\texttt{no\_tool}, \texttt{search\_books}, \texttt{search\_movies},
\texttt{search\_recipes}, \texttt{send\_email}, \texttt{translate\_text}.
The exact label map is \texttt{label\_map.json} in the repository.

\section{Reproducibility}
\label{app:repro}

Code, datasets (\texttt{data/data\_train.jsonl} / \texttt{data/data\_val.jsonl} /
\texttt{data/data\_test.jsonl}), label map, fitted temperature, and the evaluation
report are in the public repository. \texttt{requirements.txt} pins the
exact dependency versions. Reference run: \texttt{python scripts/make\_data.py}
$\to$ \texttt{python scripts/train.py} ($\sim$20\,min on a GTX 1080 Ti) $\to$
\texttt{python eval.py}. Trained weights ($\sim$2.7\,GB) are not committed;
\texttt{train.py} reproduces them.

\end{document}
